%==============================================================================
\section{Related Work}
%==============================================================================

\textbf{Threat modeling methodologies.}
STRIDE~\cite{shostack2014threat} categorizes threats into six classes (Spoofing, Tampering, Repudiation, Information Disclosure, Denial of Service, Elevation of Privilege) and remains the most widely adopted methodology. Attack trees~\cite{schneier1999attack} model adversary goals as hierarchical decompositions with AND/OR nodes, enabling quantitative risk analysis. PASTA (Process for Attack Simulation and Threat Analysis) adds a risk-centric perspective. These methodologies provide structure but require manual expertise for every system analyzed.

\textbf{Tool-assisted threat modeling.}
Microsoft Threat Modeling Tool~\cite{microsoft-tmt} generates STRIDE-based threats from data flow diagrams but requires manual diagram construction. OWASP Threat Dragon~\cite{owasp-threatdragon} offers open-source diagramming with threat rule engines. AWS Threat Composer~\cite{aws-threatcomposer} provides guided threat statement authoring with structured templates. None of these tools automate the reasoning required to generate attack trees or map threats to standardized frameworks.

\textbf{LLM-based threat analysis.}
ThreatModeling-LLM~\cite{threatmodeling-llm2024} automates threat identification for banking systems through dataset creation, prompt engineering, and model fine-tuning, achieving promising results on domain-specific threats. Bhatia et al.~\cite{canllmsthreatmodel2025} evaluate whether LLMs can threat model real-world cloud infrastructure, finding that GPT-4.1 and Gemini 2.5 Pro excel at threat identification but with varying performance across prompting strategies. ThreatLens uses retrieval-augmented generation for hardware security verification. These approaches focus on threat \emph{identification} but do not produce structured attack trees or TTP mappings.

\textbf{MITRE ATT\&CK mapping.}
Automated mapping of security text to ATT\&CK techniques has been approached through NLP classification. CVE2ATT\&CK~\cite{cve2attack2022} uses BERT-based models to map CVE descriptions to ATT\&CK techniques. SMET~\cite{smet2023} trains ATT\&CK BERT using a Siamese network to learn semantic similarity among attack actions, achieving state-of-the-art results on CVE-to-technique mapping. Kim et al.~\cite{modernbert2025} combine ModernBERT with BERTopic for tactic-level classification. These works map \emph{existing} vulnerability descriptions to techniques; \threatforest{} maps \emph{generated} attack tree steps---a distinct task requiring the model to bridge the gap between abstract attack descriptions and concrete ATT\&CK techniques.

\textbf{Domain-specific language models for cybersecurity.}
SecureBERT~\cite{aghaei2023securebert} adapts BERT to cybersecurity text, improving performance on NER and classification tasks. ATTACK-BERT~\cite{attackbert} extends sentence-transformers to produce embeddings specifically for attack action descriptions, enabling semantic similarity search against ATT\&CK technique descriptions. We leverage ATTACK-BERT as the retrieval component in our hybrid TTP mapping pipeline (\S\ref{sec:ttp}).

\textbf{Multi-agent systems.}
The Strands Agents SDK~\cite{strands2025} provides a graph-based orchestration framework for multi-agent workflows with tool use, streaming, and telemetry. While multi-agent architectures have been applied to code generation, software testing, and general reasoning, their application to security threat modeling---where outputs must be both technically accurate and structurally valid---remains unexplored.
